% !PRINTED 273
\weberpage{273}

%<<CONT
a change of sign always takes place before and after the passage
of $f_\nu$ through zero, because $f_{\nu+1}$ and $f_{\nu-1}$
have opposite signs.

But when $f(x)$ passes through zero, then by 4. a change of sign
is lost in the passage through $x$. For if $f(x)$ goes from
negative to positive values, then (cf. \S.\,32) $f'(x)$,
and hence $f_1(x)$, is positive, and the signs change thus:
\[
\begin{array}{c}
f(x), \; f_1(x) \\
- \quad + \\
+ \quad +,
\end{array}
\]
and when $f(x)$ passes from positive to negative values, then
$f'(x)$ and $f_1(x)$ are negative, so the signs are thus:
\[
\begin{array}{c}
f(x), \; f_1(x) \\
+ \quad - \\
- \quad -,
\end{array}
\]
hence in both cases a change of sign has been lost.

Now since $f_m(x)$ by 2. does not change its sign, it follows
that $V(\alpha) - V(\beta)$ is equal to the number of roots of
$f(x) = 0$ lying between $\alpha$ and $\beta$, or:

\gesperrt{The number of roots of $f(x) = 0$ between $\alpha$ and
$\beta$ is equal to the excess of the number of changes of
sign of the chain for $x = \alpha$ over the number of
changes of sign for $x = \beta$.}

If for $x = \alpha$ or $x = \beta$ one or several of the
intermediate functions of $(\mathfrak{R})$ should vanish,
then by 3. it makes no difference whether we replace this vanishing
value by a positive or by a negative one.

\begin{center}
\S.\,87.

\gesperrt{First example: spherical functions.}
\end{center}

Before we now pass on to the general methods by which
Sturm chains are to be formed, we shall consider a few examples
in which such chains present themselves through special
circumstances.
